% GENERATED FILE. Edit canonical lessons, then run npm run generate:books.
% canonical-source-hash: fnv1a64:e079931167f30786
% canonical-lessons: JA-C411-ogawa, JA-C411-izumi, JA-C411-nigoru, JA-C411-kakomu, JA-C411-sukuu

\chapter{Streams and Springs}
\label{ch:ja-streams-and-springs}
\hlchaptermodality{411}

\begin{chapteropening}
\textbf{By the end of this chapter:} \emph{I can describe water in the countryside.}

Complete the last lesson of chapter 411: I can describe water in the countryside.
\end{chapteropening}

\section[\texorpdfstring{\ja{おがわ} (ogawa)}{おがわ (ogawa)}]{\ja{おがわ} (ogawa) — a stream; a brook}

\label{lesson:JA-C411-ogawa}



\begin{quote}
Before the new one: say the Japanese for a meadow, then the Japanese for to grow thick.
\end{quote}

\subsection*{You'll want to know: \ja{おがわ}}
\textbf{\ja{おがわ}} — \emph{ogawa} — \textquotedblleft{}a stream; a brook\textquotedblright{}.

\begin{quote}
\textbf{\ja{おがわで} \ja{こどもが} \ja{あそんでいます。}} — \emph{ogawa de kodomo ga asonde imasu} — \textquotedblleft{}Children are playing in the stream.\textquotedblright{}
\end{quote}

\begin{cousinweb}
\textbf{\ja{おがわ}} (\emph{ogawa}) is a native Japanese word (wago), inherited rather than borrowed from Chinese, built from \emph{o}, \textquotedblleft{}small\textquotedblright{} and \emph{gawa}, \textquotedblleft{}river\textquotedblright{} (from \emph{kawa}); its kanji were attached later as kun'yomi (JMdict).
\end{cousinweb}

\begin{grammarlens}[title={What you've built}]
One more word for getting around.
\end{grammarlens}

\subsection*{Your turn}
\begin{itemize}
\raggedright
  \item \emph{Say it:} \emph{ogawa}
  \item \emph{Say it:} \emph{ogawa}, once more
  \item \emph{Say it:} \emph{ogawa de kodomo ga asonde imasu} — Children are playing in the stream.
  \item \emph{Recall:} say the Japanese for a meadow, then the Japanese for to grow thick, then say \emph{ogawa} again
\end{itemize}

\begin{tcolorbox}[breakable,colback=teal!4,colframe=teal!35!black,title={Before you move on}]
What does \ja{おがわ} mean? (\textquotedblleft{}A stream; a brook\textquotedblright{}.) Say it once more.
\end{tcolorbox}
\section[\texorpdfstring{\ja{いずみ} (izumi)}{いずみ (izumi)}]{\ja{いずみ} (izumi) — a spring (of water)}

\label{lesson:JA-C411-izumi}



\begin{quote}
Before the new one: say the Japanese for to grow thick, then the Japanese for a stream.
\end{quote}

\subsection*{You'll want to know: \ja{いずみ}}
\textbf{\ja{いずみ}} — \emph{izumi} — \textquotedblleft{}a spring (of water)\textquotedblright{}.

\begin{quote}
\textbf{\ja{やまの} \ja{いずみの} \ja{みずは} \ja{つめたいです。}} — \emph{yama no izumi no mizu wa tsumetai desu} — \textquotedblleft{}The water from the mountain spring is cold.\textquotedblright{}
\end{quote}

\begin{cousinweb}
\textbf{\ja{いずみ}} (\emph{izumi}) is a native Japanese word (wago), inherited rather than borrowed from Chinese; its kanji was attached later as a kun'yomi (JMdict).
\end{cousinweb}

\begin{grammarlens}[title={What you've built}]
One more word for getting around.
\end{grammarlens}

\subsection*{Your turn}
\begin{itemize}
\raggedright
  \item \emph{Say it:} \emph{izumi}
  \item \emph{Say it:} \emph{izumi}, once more
  \item \emph{Say it:} \emph{yama no izumi no mizu wa tsumetai desu} — The water from the mountain spring is cold.
  \item \emph{Recall:} say the Japanese for to grow thick, then the Japanese for a stream, then say \emph{izumi} again
\end{itemize}

\begin{tcolorbox}[breakable,colback=teal!4,colframe=teal!35!black,title={Before you move on}]
What does \ja{いずみ} mean? (\textquotedblleft{}A spring (of water)\textquotedblright{}.) Say it once more.
\end{tcolorbox}
\section[\texorpdfstring{\ja{にごる} (nigoru)}{にごる (nigoru)}]{\ja{にごる} (nigoru) — to become muddy; to cloud over (of water)}

\label{lesson:JA-C411-nigoru}



\begin{quote}
Before the new one: say the Japanese for a stream, then the Japanese for a spring (of water).
\end{quote}

\subsection*{You'll want to know: \ja{にごる}}
\textbf{\ja{にごる}} — \emph{nigoru} — \textquotedblleft{}to become muddy; to cloud over (of water)\textquotedblright{}.

\begin{quote}
\textbf{\ja{あめで} \ja{かわの} \ja{みずが} \ja{にごりました。}} — \emph{ame de kawa no mizu ga nigorimashita} — \textquotedblleft{}The rain made the river water muddy.\textquotedblright{}
\end{quote}

Polite \textbf{\ja{にごります}}, \emph{nigorimasu}; a godan verb. Past (polite): \textbf{\ja{かわの} \ja{みずが} \ja{にごりました。}}, \emph{kawa no mizu ga nigorimashita}.

\begin{cousinweb}
\textbf{\ja{にごる}} (\emph{nigoru}) is a native Japanese word (wago), inherited rather than borrowed from Chinese; its kanji was attached later as a kun'yomi (JMdict).
\end{cousinweb}

\begin{grammarlens}[title={What you've built}]
One more word for getting around.
\end{grammarlens}

\subsection*{Your turn}
\begin{itemize}
\raggedright
  \item \emph{Say it:} \emph{nigoru}
  \item \emph{Say it:} \emph{nigoru}, once more
  \item \emph{Say it:} \emph{ame de kawa no mizu ga nigorimashita} — The rain made the river water muddy.
  \item \emph{Recall:} say the Japanese for a stream, then the Japanese for a spring (of water), then say \emph{nigoru} again
\end{itemize}

\begin{tcolorbox}[breakable,colback=teal!4,colframe=teal!35!black,title={Before you move on}]
What does \ja{にごる} mean? (\textquotedblleft{}To become muddy; to cloud over (of water)\textquotedblright{}.) Say it once more.
\end{tcolorbox}
\section[\texorpdfstring{\ja{かこむ} (kakomu)}{かこむ (kakomu)}]{\ja{かこむ} (kakomu) — to surround}

\label{lesson:JA-C411-kakomu}



\begin{quote}
Before the new one: say the Japanese for a spring (of water), then the Japanese for to become muddy.
\end{quote}

\subsection*{You'll want to know: \ja{かこむ}}
\textbf{\ja{かこむ}} — \emph{kakomu} — \textquotedblleft{}to surround\textquotedblright{}.

\begin{quote}
\textbf{\ja{もりが} \ja{いずみを} \ja{かこんでいます。}} — \emph{mori ga izumi o kakonde imasu} — \textquotedblleft{}Forest surrounds the spring.\textquotedblright{}
\end{quote}

Polite \textbf{\ja{かこみます}}, \emph{kakomimasu}; a godan verb. Past (polite): \textbf{\ja{みんなが} \ja{せんせいを} \ja{かこみました。}}, \emph{minna ga sensei o kakomimashita}.

\begin{cousinweb}
\textbf{\ja{かこむ}} (\emph{kakomu}) is a native Japanese word (wago), inherited rather than borrowed from Chinese; its kanji was attached later as a kun'yomi (JMdict).
\end{cousinweb}

\begin{grammarlens}[title={What you've built}]
One more word for getting around.
\end{grammarlens}

\subsection*{Your turn}
\begin{itemize}
\raggedright
  \item \emph{Say it:} \emph{kakomu}
  \item \emph{Say it:} \emph{kakomu}, once more
  \item \emph{Say it:} \emph{mori ga izumi o kakonde imasu} — Forest surrounds the spring.
  \item \emph{Recall:} say the Japanese for a spring (of water), then the Japanese for to become muddy, then say \emph{kakomu} again
\end{itemize}

\begin{tcolorbox}[breakable,colback=teal!4,colframe=teal!35!black,title={Before you move on}]
What does \ja{かこむ} mean? (\textquotedblleft{}To surround\textquotedblright{}.) Say it once more.
\end{tcolorbox}
\section[\texorpdfstring{\ja{すくう} (sukuu)}{すくう (sukuu)}]{\ja{すくう} (sukuu) — to scoop up}

\label{lesson:JA-C411-sukuu}



\begin{quote}
Before the new one: say the Japanese for to become muddy, then the Japanese for to surround.
\end{quote}

\subsection*{You'll want to know: \ja{すくう}}
\textbf{\ja{すくう}} — \emph{sukuu} — \textquotedblleft{}to scoop up\textquotedblright{}.

\begin{quote}
\textbf{\ja{いずみの} \ja{みずを} \ja{てで} \ja{すくいました。}} — \emph{izumi no mizu o te de sukuimashita} — \textquotedblleft{}I scooped up water from the spring with my hands.\textquotedblright{}
\end{quote}

Polite \textbf{\ja{すくいます}}, \emph{sukuimasu}; a godan verb. Past (polite): \textbf{\ja{こどもが} \ja{みずを} \ja{すくいました。}}, \emph{kodomo ga mizu o sukuimashita}.

\begin{cousinweb}
\textbf{\ja{すくう}} (\emph{sukuu}) is a native Japanese word (wago), inherited rather than borrowed from Chinese; its kanji was attached later as a kun'yomi (JMdict).
\end{cousinweb}

\begin{grammarlens}[title={What you've built}]
One more word for getting around.
\end{grammarlens}

\subsection*{Your turn}
\begin{itemize}
\raggedright
  \item \emph{Say it:} \emph{sukuu}
  \item \emph{Say it:} \emph{sukuu}, once more
  \item \emph{Say it:} \emph{izumi no mizu o te de sukuimashita} — I scooped up water from the spring with my hands.
  \item \emph{Recall:} say the Japanese for to become muddy, then the Japanese for to surround, then say \emph{sukuu} again
\end{itemize}

\begin{tcolorbox}[breakable,colback=teal!4,colframe=teal!35!black,title={Before you move on}]
What does \ja{すくう} mean? (\textquotedblleft{}To scoop up\textquotedblright{}.) Say it once more.
\end{tcolorbox}
